% fig_vl01_planck.tex — Planck gegen Rayleigh-Jeans: das Spektrum, an dem
% die klassische Physik scheitert, und der Grund dafür.
%
% Folgt den Konventionen von fig_vl11_kasten.tex.
%
% Was die Abbildung zeigt
%   (a) die Energiedichte U(nu) für drei Temperaturen, daneben die
%       klassische Vorhersage, die oben aus dem Bild herausläuft;
%   (b) warum: die mittlere Energie einer Mode, klassisch kT, tatsächlich
%       hnu/(exp(hnu/kT)-1) — oberhalb von hnu ~ kT friert sie ein.
%   Die beiden Panels hängen multiplikativ zusammen,
%       U(nu) = (8 pi nu^2/c^3) * <E>(nu,T),
%   also (a) = Modendichte x (b).  Genau der zweite Faktor ist es, den die
%   klassische Physik falsch bekommt; die Modendichte ist klassisch richtig.
%   Panel (b) ist zugleich das Verhältnis der beiden Kurven aus (a):
%       U/U_RJ = x/(e^x - 1)  mit  x = h nu/kT.
%
% Geometrie aus der Formel, nichts freihand
%   Planck (die Formel der Vorlesung):
%       U(nu) = (8 pi h nu^3/c^3) / (exp(h nu/kT) - 1)
%   Rayleigh-Jeans ist ihr Grenzfall h nu << kT:
%       U_RJ(nu) = 8 pi nu^2 k T / c^3
%   Mit nu = s * 1e14 Hz und U in 1e-15 J m^-3 Hz^-1 wird daraus
%       8 pi h/c^3 = 6.1806e-58  ->  U    = 0.61806 * s^3/(exp(a_T s) - 1),
%       8 pi k/c^3 = 1.28786e-47 ->  U_RJ = 1.28786e-4 * T * s^2,
%   mit a_T = h*1e14/(kT) = 4799.2/T, also a = 1.19980, 0.95984, 0.79987
%   für T = 4000, 5000, 6000 K.  Damit liegen die Maxima von selbst
%   richtig, wo Wien sie verlangt (h nu_max = 2.8214 kT):
%       s_max = 2.35, 2.94, 3.53  mit  U_max = 0.51, 0.99, 1.72,
%   d.h. die Lage wandert wie T und die Höhe wächst wie T^3.  Beides liest
%   man im Bild ab, ohne dass es hingeschrieben werden müsste.
%   Panel (b) ist dieselbe Planck-Formel in der dimensionslosen Variablen
%   x = h nu/kT:  <E>/kT = x/(exp(x)-1),  klassisch = 1.
%
% Entwurfsentscheidungen
%   * Farbe trägt die Temperatur, kalt -> heiß (spacecadet, cgblue,
%     birrenorange), und sonst nichts.  Beide klassischen Kurven —
%     Rayleigh-Jeans in (a) und die Gleichverteilungslinie <E> = kT in (b)
%     — sind neutral grau gestrichelt: gescheiterte Referenz, keine Daten.
%     Die eine Kurve in (b) ist temperaturunabhängig, also bekommt sie mit
%     birrengreen bewusst eine Farbe, die im Temperaturverlauf von (a)
%     nicht vorkommt.
%   * Rayleigh-Jeans ist für die *kälteste* der drei Temperaturen
%     gezeichnet, also für die flachste der drei klassischen Parabeln.
%     Sie verlässt den Rahmen trotzdem, bevor das zugehörige
%     Planck-Spektrum überhaupt sein Maximum erreicht, und übertrifft bei
%     nu = 2e14 Hz bereits das 6000-K-Spektrum: die Katastrophe hängt
%     nicht an der Temperatur.  Bei kleinen nu liegt die graue Kurve exakt
%     auf der navy Planck-Kurve — klassisch ist eben nicht alles falsch,
%     und genau das sagt (b) noch einmal quantitativ.
%   * Die Lücke zwischen kT und <E> in (b) ist schraffiert, nicht gefüllt
%     (Graustufendruck).  Sie ist als TikZ-Pfad gezeichnet und nicht mit
%     `fill between`: diese Bibliothek zerstört hier die Bounding Box der
%     tikzpicture, das rechte Panel wird dann abgeschnitten.
%   * Echte Einheiten (1e14 Hz, 1e-15 J m^-3 Hz^-1, Kelvin), damit die
%     Größenordnung mitkommt; 6000 K liegt bei der Sonnenoberfläche.
%   * Beschriftung direkt am Objekt statt in einer Legende.  Die
%     Achsenbeschriftungen stehen mittig unter bzw. neben der Achse statt
%     am Achsenende (tpcplot-Voreinstellung): am Achsenende würde die von
%     (a) in das rechte Panel hineinlaufen.

\begin{center}
\begin{tikzpicture}

% ─────────────────────────────────────────────────────────────
%  (a)  Das Spektrum: Planck für drei T, Rayleigh-Jeans daneben
% ─────────────────────────────────────────────────────────────
\begin{axis}[
    name=spectrum,
    tpcplot,
    width=8.2cm, height=5.8cm,
    xmin=0, xmax=8.8, ymin=0, ymax=2.6,
    xtick={0,2,4,6,8}, ytick={0,0.5,1,1.5,2,2.5},
    xlabel={$\nu\;[10^{14}\,\mathrm{Hz}]$},
    ylabel={$U(\nu)\;[10^{-15}\,\mathrm{J\,m^{-3}\,Hz^{-1}}]$},
    every axis x label/.style={at={(axis description cs:0.5,-0.13)},
                               anchor=north, font=\small},
    every axis y label/.style={at={(axis description cs:-0.135,0.5)},
                               rotate=90, anchor=south, font=\scriptsize},
    title={(a) \tpcifenglish{The spectrum}{Das Spektrum}},
]

  % Planck, U = 0.61806 s^3/(exp(a_T s) - 1)
  \addplot[spacecadet,   line width=1.1pt, domain=0.01:8.8, samples=220]
    {0.61806*x^3/(exp(1.199800*x)-1)};
  \addplot[cgblue,       line width=1.1pt, domain=0.01:8.8, samples=220]
    {0.61806*x^3/(exp(0.959840*x)-1)};
  \addplot[birrenorange, line width=1.1pt, domain=0.01:8.8, samples=220]
    {0.61806*x^3/(exp(0.799867*x)-1)};

  % Rayleigh-Jeans für T = 4000 K, U_RJ = 8 pi nu^2 kT/c^3: der Grenzfall
  % h nu << kT der Planck-Formel.  Grau, weil klassische Referenz.  Sie
  % wird zuletzt gezeichnet, damit man sie bis nu ~ 1e14 Hz genau auf der
  % navy Planck-Kurve liegen sieht, bevor sie sich löst und oben aus dem
  % Rahmen läuft.
  \addplot[birrengray, line width=0.9pt, densely dashed,
           domain=0:2.30, samples=60] {0.515144*x^2};

  % Temperaturen direkt an ihren Maxima (h nu_max = 2.8214 kT)
  \node[lbl, font=\scriptsize, text=birrenorange, anchor=south]
    at (axis cs:3.53,1.76) {$6000\,\mathrm{K}$};
  \node[lbl, font=\scriptsize, text=cgblue, anchor=south]
    at (axis cs:2.94,1.03) {$5000\,\mathrm{K}$};
  \node[lbl, font=\scriptsize, text=spacecadet, anchor=south]
    at (axis cs:2.35,0.55) {$4000\,\mathrm{K}$};

  % Die klassische Kurve und ihr Schicksal
  \node[lbl, font=\scriptsize, text=birrengray, anchor=north west,
        align=left, inner sep=1pt]
    at (axis cs:2.45,2.58)
    {Rayleigh--Jeans ($4000\,\mathrm{K}$)\\
     $U\propto\nu^{2}$: \tpcifenglish{UV catastrophe}{UV-Katastrophe}};
\end{axis}

% ─────────────────────────────────────────────────────────────
%  (b)  Warum: die mittlere Energie einer Mode
% ─────────────────────────────────────────────────────────────
\begin{axis}[
    tpcplot,
    at={(spectrum.south east)}, anchor=south west, xshift=10mm,
    width=5.4cm, height=5.8cm,
    xmin=0, xmax=6, ymin=0, ymax=1.28,
    xtick={0,1,2,3,4,5,6}, ytick={0,0.5,1},
    xlabel={$h\nu/kT$},
    ylabel={$\langle E\rangle/kT$},
    every axis x label/.style={at={(axis description cs:0.5,-0.13)},
                               anchor=north, font=\small},
    every axis y label/.style={at={(axis description cs:-0.16,0.5)},
                               rotate=90, anchor=south, font=\small},
    title={(b) \tpcifenglish{Energy per mode}{Energie pro Mode}},
]

  % Die Lücke zwischen klassischer Gleichverteilung und Planck:
  % schraffiert, nicht gefüllt.  Als Pfad, nicht mit `fill between`.
  \draw[forbidden, draw=none]
    (axis cs:0.02,1) -- plot[domain=0.02:6, samples=140, variable=\t]
      (axis cs:{\t},{\t/(exp(\t)-1)}) -- (axis cs:6,1) -- cycle;

  % Klassisch: jede Mode trägt kT, unabhängig von nu.
  \addplot[levellinegrey, domain=0:6, samples=2] {1};
  \node[lbl, font=\scriptsize, text=birrengray, anchor=south west,
        inner sep=1pt]
    at (axis cs:0.12,1.04)
    {\tpcifenglish{classical}{klassisch}: $\langle E\rangle=kT$};

  % Planck: <E>/kT = x/(exp(x)-1)
  \addplot[birrengreen, line width=1.2pt, domain=0.02:6, samples=200]
    {x/(exp(x)-1)};
  \node[lbl, font=\scriptsize, text=birrengreen, anchor=center,
        fill=white, fill opacity=0.92, text opacity=1, inner sep=1.5pt]
    at (axis cs:3.45,0.63)
    {$\langle E\rangle=\dfrac{h\nu}{e^{h\nu/kT}-1}$};

  % Was die Schraffur bedeutet
  \node[lbl, font=\scriptsize, text=birrengray, anchor=east,
        fill=white, fill opacity=0.92, text opacity=1, inner sep=1.5pt]
    at (axis cs:5.9,0.28)
    {\tpcifenglish{frozen out}{eingefroren}};
\end{axis}

\end{tikzpicture}
\end{center}
